% ==============================================================================
% BLOQUE 2: ANATOMÍA CCC, PRUEBAS DE ACEPTACIÓN (BDD), DoD, GRANULARIDAD E INVEST
% Diapositivas 26 a 57 del Seminario 2
% ==============================================================================

\chapter[Anatomía «CCC», Pruebas BDD, Granularidad e INVEST]{Anatomía «CCC», Pruebas de Aceptación,\\ Granularidad y Modelo INVEST}

\cuadrofuentes{
\textbf{Alcance del Bloque:} Diapositivas 26--57 del Seminario 2.\\
\textbf{Conceptos Clave:} Fórmula \textbf{CCC} de Ron Jeffries (\textit{Card, Conversation, Confirmation}), anatomía física de la tarjeta (anverso y reverso), búsqueda de detalles en la Conversación, Pruebas de Aceptación de caja negra, formato BDD (\texttt{Dado / Cuando / Entonces}), Definición de Terminado (\textit{Definition of Done}), jerarquía de granularidad (\textit{Temas, Épicas, HU, Tareas}), gestión de Requisitos No Funcionales (RNF) y modelo \textbf{I.N.V.E.S.T.} de Bill Wake.
}

\section{Elementos de una Historia de Usuario: El Modelo ``CCC'' (Ron Jeffries)}

En el año 2001, \textbf{Ron Jeffries} (uno de los tres fundadores de \textit{Extreme Programming}) formuló la regla de las \textbf{Tres ``C''} para explicar que una Historia de Usuario es mucho más que el texto escrito en una cartulina:

\begin{center}
\resizebox{0.92\linewidth}{!}{%
\begin{tikzpicture}[
    cbox/.style={rectangle, rounded corners=5pt, draw=mdaPrimary, thick, text width=4.2cm, minimum width=4.5cm, minimum height=3.3cm, align=center, font=\sffamily},
    arrow/.style={-{Stealth[length=3.5mm, width=2.5mm]}, ultra thick, color=mdaAccent}
]
    \node[cbox, fill=mdaYellowCard, draw=mdaYellowBorder] (c1) at (0,0) {
        {\Large\bfseries\color{mdaDarkBlue}1. Tarjeta}\\[2pt]
        {\normalsize\bfseries\color{mdaOrange}(Card)}\\[6pt]
        {\scriptsize 2 ó 3 frases que describen}\\[-2pt]
        {\scriptsize la intención de la HU.}\\[3pt]
        {\scriptsize\itshape Es una promesa; si tiene}\\[-2pt]
        {\scriptsize\itshape demasiados detalles, limita}\\[-2pt]
        {\scriptsize\itshape la conversación con el cliente.}
    };

    \node[cbox, fill=mdaLightAccent, draw=mdaPrimary, right=0.6cm of c1] (c2) {
        {\Large\bfseries\color{mdaDarkBlue}2. Conversación}\\[2pt]
        {\normalsize\bfseries\color{mdaPrimary}(Conversation)}\\[6pt]
        {\scriptsize El verdadero trabajo ágil.}\\[-2pt]
        {\scriptsize Detalles obtenidos del}\\[-2pt]
        {\scriptsize diálogo continuo entre el}\\[-2pt]
        {\scriptsize equipo de desarrollo y los}\\[-2pt]
        {\scriptsize usuarios o clientes.}
    };

    \node[cbox, fill=mdaLightGreen, draw=mdaGreen, right=0.6cm of c2] (c3) {
        {\Large\bfseries\color{mdaDarkBlue}3. Confirmación}\\[2pt]
        {\normalsize\bfseries\color{mdaGreen}(Confirmation)}\\[6pt]
        {\scriptsize Pruebas de aceptación.}\\[-2pt]
        {\scriptsize Cómo el equipo y el cliente}\\[-2pt]
        {\scriptsize confirmarán sin ambigüedades}\\[-2pt]
        {\scriptsize que la HU ha sido}\\[-2pt]
        {\scriptsize implementada y terminada.}
    };

    \draw[arrow] (c1) -- (c2);
    \draw[arrow] (c2) -- (c3);
\end{tikzpicture}%
}
\end{center}

\subsection{Información y Narración en la Tarjeta (Diapositivas 28--30)}
\begin{itemize}[itemsep=3pt]
    \item \textbf{Descripción corta sin detalles exhaustivos:} La tarjeta presenta una descripción concisa junto con pequeñas notas iniciales (p. ej., \textit{``Un cliente puede pagar un viaje usando una tarjeta de crédito. Nota: Se pueden aceptar Visa, MasterCard y American Express; discutir si se aceptan otras más''}).
    \item \textbf{Peligro del exceso de detalle:} \textbf{Una tarjeta con demasiados detalles limita la conversación con el cliente}, induciendo a los programadores a asumir que todo está decidido y bloqueando soluciones más creativas o simples.
    \item \textbf{Delimitación del alcance:} La parte exacta a abordar por la HU se delimita en la conversación y se plasma usando las \textbf{pruebas de aceptación}, las cuales recogen:
    \begin{enumerate}
        \item Flujos alternativos en la actividad.
        \item Límites de aceptación (condiciones de contorno y restricciones).
        \item Pruebas de aceptación verificables por parte de los clientes.
    \end{enumerate}
\end{itemize}

\section{Anatomía Canónica de la Tarjeta Física de HU (Diapositivas 31--34)}

En la diapositiva 31 se establece la estructura oficial de una tarjeta de Historia de Usuario dividida en \textbf{Anverso} (identificación, descripción, estimación, prioridad y dependencias) y \textbf{Reverso} (pruebas de aceptación):

\begin{center}
\resizebox{0.92\linewidth}{!}{%
\begin{tikzpicture}[
    card/.style={rectangle, draw=mdaYellowBorder!90!black, fill=mdaYellowCard, very thick, minimum width=10.2cm, align=left, font=\sffamily\small},
    tag/.style={rectangle, rounded corners=6pt, draw=mdaDarkBlue, fill=white, thick, font=\sffamily\scriptsize\bfseries, inner sep=4pt, text=mdaDarkBlue},
    line/.style={-{Circle[length=1.8mm]}, thick, color=mdaDarkBlue}
]
    % ANVERSO
    \node[font=\sffamily\bfseries\color{mdaDarkBlue}] at (0, 2.5) {ANVERSO DE LA TARJETA};
    \node[card, minimum height=4.0cm] (anverso) at (0, 0.2) {};
    \draw[mdaYellowBorder!90!black, thick] (-5.1, 1.45) -- (5.1, 1.45);
    
    \node[anchor=west, font=\sffamily\bfseries\color{mdaDarkBlue}] (idval) at (-4.8, 1.8) {5};
    \node[ font=\sffamily\bfseries\color{mdaDarkBlue}] (titval) at (0, 1.8) {Préstamo de Libro};
    
    \node[anchor=north west, text width=9.6cm, font=\sffamily\footnotesize] (descval) at (-4.8, 1.25) {
        Cómo cliente quiero que los socios puedan pedir prestado un libro, indicando su número de socio y la referencia del libro, siempre y cuando no tengan ya tres libros en préstamo en ese momento.
    };
    
    \node[anchor=west, font=\sffamily\footnotesize] (estval) at (-4.8, -0.75) {Estimación: \textbf{4}};
    \node[anchor=west, font=\sffamily\footnotesize] (prival) at (-4.8, -1.35) {Prioridad: \textbf{300}};
    \node[anchor=east, font=\sffamily\footnotesize] (depval) at (4.8, -1.35) {Dependiente de: \textbf{1, 2}};

    % Etiquetas Anverso
    \node[tag] (tagid) at (-6.4, 1.8) {ID};
    \node[tag] (tagtit) at (6.4, 1.8) {TÍTULO};
    \node[tag] (tagdesc) at (6.5, 0.4) {DESCRIPCIÓN};
    \node[tag] (tagest) at (-6.5, -0.75) {ESTIMACIÓN};
    \node[tag] (tagpri) at (-6.5, -1.35) {PRIORIDAD};
    \node[tag] (tagdep) at (6.6, -1.35) {DEPENDENCIAS};

    \draw[line] (tagid.east) -- (-4.85, 1.8);
    \draw[line] (tagtit.west) -- (1.8, 1.8);
    \draw[line] (tagdesc.west) -- (4.2, 0.4);
    \draw[line] (tagest.east) -- (-4.85, -0.75);
    \draw[line] (tagpri.east) -- (-4.85, -1.35);
    \draw[line] (tagdep.west) -- (4.75, -1.35);

    % REVERSO
    \node[font=\sffamily\bfseries\color{mdaDarkBlue}] at (0, -2.5) {REVERSO DE LA TARJETA};
    \node[card, minimum height=4.4cm] (reverso) at (0, -5.0) {};
    \node[anchor=north west, text width=9.6cm, font=\sffamily\footnotesize] at (-4.8, -3.1) {
        $\bullet$ Introducir un número de socio incorrecto y comprobar que se indica error.\\[3pt]
        $\bullet$ Introducir un socio que ya tiene 3 libros en préstamo y comprobar que se indica error.\\[3pt]
        $\bullet$ Introducir un libro del que no hay ejemplares y comprobar que se indica un error.\\[3pt]
        $\bullet$ Introducir todos los datos correctos y comprobar que el número de ejemplares disponibles del libro disminuye y el número de préstamos del socio aumenta en uno.
    };
    \node[tag, align=center] (tagrev) at (6.6, -5.0) {PRUEBAS DE\\ACEPTACIÓN};
    \draw[line] (tagrev.west) -- (4.4, -5.0);
\end{tikzpicture}%
}
\end{center}

En herramientas digitales como \textbf{Jira} (diapositiva 34) o plantillas corporativas (diapositivas 32--33), estos mismos campos se representan mediante atributos equivalentes: identificador único, título, narrativa \texttt{Como/Quiero/Para}, condiciones de completitud, prioridad numérica, coste/puntos estimados, iteración asignada y persona encargada.

\section{Buscar Detalles de la HU: La Conversación (Diapositivas 35--36)}

\begin{aviso}[Principio Clave de la Conversación (Diapositiva 35)]
\textbf{``La tarjeta no contiene estas respuestas; su propósito es PROVOCAR estas preguntas.''}
\end{aviso}

Por ejemplo, ante la tarjeta mínima \textit{``Como cliente, quiero cancelar una reserva de hotel''}, la etapa de \textbf{Conversación} entre desarrolladores y usuarios debe hacer aflorar preguntas como:
\begin{itemize}[itemsep=2pt]
    \item \textit{¿Qué coste de penalización tiene la cancelación?}
    \item \textit{¿Todos los niveles de clientes tienen el mismo coste o varía si es cliente Premium?}
    \item \textit{¿Se puede cancelar en cualquier momento o hay límite de horas previo?}
    \item \textit{¿Se permite una cancelación total y también parcial?}
    \item \textit{¿A quién (departamento/hotel/cliente) se le notifica automáticamente?}
\end{itemize}

\textbf{¿Por qué son vitales las conversaciones? (Diapositiva 36):}
\begin{enumerate}[label=\textbf{\arabic*.}, itemsep=2pt]
    \item Ayudan al equipo a comprender las necesidades y deseos reales del usuario y a producir \textbf{soluciones creativas} a los problemas.
    \item Favorecen la \textbf{solidez de las relaciones} de confianza entre los miembros del equipo y los usuarios.
    \item Deben reflejar la \textbf{voz de todos los interesados} (\textit{stakeholders}) en el producto.
\end{enumerate}

\section{Pruebas de Aceptación: La Confirmación y Formato BDD (Diapositivas 37--43)}

\begin{definicion}[Criterios o Pruebas de Aceptación]
Los criterios o pruebas de aceptación son las \textbf{condiciones objetivas que una Historia de Usuario debe cumplir para considerarse completada} cuando es desarrollada y desplegada por el equipo. Funcionan como una guía que asegura que lo que se construye es realmente lo que el negocio necesita.
\end{definicion}

\subsection{Reglas y Propiedades de las Pruebas de Aceptación}
\begin{itemize}[itemsep=3pt]
    \item \textbf{Claridad y precisión:} No deben dejar espacio a interpretaciones subjetivas ni ambigüedades.
    \item \textbf{Independencia de la implementación:} No describen \textit{cómo} programarlo internamente, sino \textbf{\textit{qué} resultado observable se espera}.
    \item \textbf{Lenguaje de negocio:} Se escriben siempre en el lenguaje del ámbito de negocios.
    \item \textbf{Propiedad del Product Owner:} Aunque cualquiera en el equipo puede escribir una prueba de aceptación, \textbf{el Product Owner es el principal propietario de la prueba}.
    \item \textbf{Pruebas de Caja Negra (\textit{Black-box testing}):} Solo verifican que las salidas del sistema cumplen las condiciones de satisfacción indicadas ante determinadas entradas, sin preocuparse de cómo se consiguen internamente los resultados.
    \item \textbf{Sincronía de iteración:} Se implementan y verifican \textbf{en el transcurso de la misma iteración} en la que se implementa la propia Historia de Usuario.
    \item \textbf{Complementariedad con Pruebas Unitarias:} El equipo de desarrollo realiza \textbf{pruebas de unidad} (caja blanca) sobre los módulos de código desarrollados, \textit{independientemente} de las pruebas de aceptación definidas para la HU.
\end{itemize}

\begin{ejemplo}[Criterios de Aceptación Clásicos y Formato BDD (Diapositivas 40--43)]
\textbf{Caso 1 (Lista de comprobación -- Cancelar reserva previa):}
\begin{itemize}[itemsep=1pt]
    \item Comprobar que un \textbf{Cliente Premium} puede cancelar en el mismo día sin coste adicional.
    \item Comprobar que a un \textbf{Cliente No-Premium} se le carga un $10\%$ al cancelar el mismo día.
    \item Confirmar que se envía un correo electrónico de confirmación al cliente.
    \item Confirmar que el hotel recibe notificación automática de la cancelación.
\end{itemize}

\textbf{Caso 2 (Lista de comprobación -- Pagar viaje con tarjeta de crédito):}
\begin{itemize}[itemsep=1pt]
    \item Funciona correctamente con tarjetas \textit{Visa, MasterCard y American Express}.
    \item Falla (no permitida) con tarjeta \textit{Diner's Club}.
    \item Rechaza adecuadamente una tarjeta cuya fecha haya expirado.
    \item Maneja correctamente diferentes montos de cargo y falla ante un número de tarjeta inválido o vacío.
\end{itemize}

\textbf{Caso 3 (Formato BDD -- Behavior-Driven Development: \texttt{Dado / Cuando / Entonces}):}
\begin{center}
\small
\renewcommand{\arraystretch}{1.25}
\begin{tabular}{r l}
\textbf{\color{mdaGreen}Dado} & \textit{[contexto inicial / estado previo del sistema]} \\
$\boldsymbol{\longrightarrow}$ \textbf{\color{mdaGreen}Cuando} & \textit{[acción realizada por el usuario o evento desencadenante]} \\
$\boldsymbol{\longrightarrow}$ \textbf{\color{mdaGreen}Entonces} & \textit{[resultado esperado / respuesta observable del sistema]}
\end{tabular}
\end{center}
\begin{lstlisting}
DADO un sistema donde existe un usuario con nombre "idprueba"
CUANDO alguien intenta registrarse con nombre "idprueba"
ENTONCES el sistema muestra un mensaje indicando que el nombre "idprueba" ya existe.
\end{lstlisting}
\end{ejemplo}

\section{Características de las HU, ``HU Terminada'' (Done) y ¿Por qué usar HU?}

\subsection{Síntesis de Características y Ventajas de las HU (Diapositivas 44 y 48)}
\begin{itemize}[itemsep=3pt]
    \item Son \textbf{cortas y fáciles de leer}, entendibles por desarrolladores, interesados y usuarios.
    \item Representan \textbf{incrementos pequeños de funcionalidad valorada}, desarrollables en un período de días o pocas semanas.
    \item Son \textbf{relativamente fáciles de estimar} porque el esfuerzo de implementación puede determinarse rápidamente.
    \item Se gestionan en \textbf{listas organizadas dinámicas} (no en documentos pesados) que pueden reordenarse a medida que emerge nueva información (\textbf{desarrollo oportunista / requisitos emergentes}).
    \item Necesitan \textbf{poco o ningún mantenimiento} y se pueden desechar con seguridad después de la implementación: las HU y el propio código creado sirven como documentación incremental.
    \item Permiten mantener una \textbf{relación cercana con el cliente} y aplazar los detalles para cuando sean realmente necesarios.
\end{itemize}

\subsection{HU Terminada: \textit{Definition of Done} (DoD) (Diapositivas 45--47)}

\begin{definicion}[Definición de ``Terminada'' (Definition of Done -- DoD)]
La definición de ``terminado'' (\textit{Done}) es un \textbf{acuerdo formal de calidad entre los desarrolladores y los clientes / Product Owner}. Se puede establecer una definición de ``terminada'' global para todas las HU del proyecto o un ``terminado'' específico adicional para determinadas HU.
\end{definicion}

\textbf{ los 6 aspectos esenciales a considerar en la definición de ``terminada'' (Diapositiva 47):}
\begin{enumerate}[label=\textbf{\arabic*.}, itemsep=2pt]
    \item \textbf{Tipos y grado de prueba:} Superación de pruebas de unidad y pruebas de integración.
    \item \textbf{Calidad del código:} Código documentado, revisado por pares (\textit{code review}) y sin incidencias críticas.
    \item \textbf{Entorno de despliegue donde concluye el trabajo:} En preproducción (\textit{prepro}), \textit{testing} o producción.
    \item \textbf{Nivel de documentación} técnica y de usuario entregada.
    \item \textbf{Nivel de calidad del producto software:} Cumplimiento de métricas de análisis estático, cobertura, etc.
    \item \textbf{Criterios de aceptación comprobados y aprobados} formalmente por el Product Owner.
\end{enumerate}

\section{Granularidad de las HU y Requisitos No Funcionales (Diapositivas 49--54)}

\subsection{Jerarquía de Granularidad: Temas, Épicas, HU y Tareas}

En un proyecto ágil coexisten elementos de distinto tamaño y nivel de abstracción. A medida que una necesidad se acerca a la parte alta de la pila para ser implementada en la siguiente iteración, se descompone desde conceptos amplios hacia unidades más pequeñas:

\begin{center}
\resizebox{0.95\linewidth}{!}{%
\begin{tikzpicture}[
    level/.style={trapezium, trapezium angle=72, draw=mdaDarkBlue, thick, align=center, font=\sffamily\footnotesize, minimum height=1.05cm},
    desc/.style={rectangle, rounded corners=3pt, draw=mdaBorderGray, fill=mdaGray, thick, text width=8.2cm, minimum height=1.05cm, align=left, font=\sffamily\scriptsize, inner sep=5pt}
]
    \node[level, fill=mdaDarkBlue, text=white, minimum width=3.2cm, trapezium stretches=true] (l1) at (0, 3.6) {\textbf{1. TEMAS}\\[-1pt](Themes)};
    \node[level, fill=mdaPrimary, text=white, minimum width=4.4cm, trapezium stretches=true] (l2) at (0, 2.4) {\textbf{2. ÉPICAS}\\[-1pt](Epics)};
    \node[level, fill=mdaAccent, text=white, minimum width=5.6cm, trapezium stretches=true] (l3) at (0, 1.2) {\textbf{3. HISTORIAS DE USUARIO}\\[-1pt](User Stories)};
    \node[level, fill=mdaTeal, text=white, minimum width=6.8cm, trapezium stretches=true] (l4) at (0, 0.0) {\textbf{4. TAREAS TÉCNICAS}\\[-1pt](Tasks)};

    \node[desc] at (6.8, 3.6) {\textbf{Grandes proyectos o subproductos}, peticiones globales sin más análisis ni detalles. \textit{Ej: Planificador de viajes, gestor de ofertas de hoteles, sitio Web para socios.}};
    \node[desc] at (6.8, 2.4) {\textbf{Súper historias de usuario}, más concretas que los Temas. Conjunto de HU relacionadas entre sí. \textit{Ej: Sistema de búsqueda para hoteles, filtros aplicados a búsquedas, comparador avanzado.}};
    \node[desc] at (6.8, 1.2) {\textbf{Descripción simple de una tarea concisa que aporta valor} en 1 iteración. \textit{Ej: Como cliente quiero buscar un hotel en una localidad determinada; como socio quiero pagar con tarjeta.}};
    \node[desc] at (6.8, 0.0) {\textbf{División de una HU por necesidades técnicas o de desarrollo} ($0.5$ a $3$--$4$ días). \textit{Ej: Crear pantalla de resultados, crear métodos de consulta a BBDD, probar integración del módulo.}};
\end{tikzpicture}%
}
\end{center}

\subsection{Requisitos No Funcionales (RNF) en Ágil (Diapositivas 51--54)}

Los \textbf{Requisitos No Funcionales (RNF)} describen cualidades o restricciones del sistema (rendimiento, disponibilidad, seguridad, usabilidad) que no se relacionan de forma directa con un comportamiento funcional concreto, pudiendo aplicarse \textbf{al sistema en su conjunto} o \textbf{a determinadas funciones específicas}.

\begin{teorema}[Las 4+1 Formas de Tratar los Requisitos No Funcionales (Diapositiva 52)]
En un proyecto ágil basado en Historias de Usuario, los RNF pueden tratarse como:
\begin{enumerate}[label=\textbf{\arabic*.}, itemsep=3pt]
    \item \textbf{Restricciones asignadas a cualquiera de los elementos de las pilas} (Temas, Épicas, Historias de Usuario o Tareas) o redactadas directamente con formato de HU.
    \item \textbf{Criterios de aceptación} asociados a una Historia de Usuario funcional.
    \item \textbf{Parte de la definición de terminada (\textit{Definition of Done})} cuando afectan globalmente a todas las historias.
    \item \textbf{Documento de especificación específico} para los requisitos no funcionales globales al sistema.
\end{enumerate}
\end{teorema}

\begin{ejemplo}[Ejemplos Oficiales de Requisitos No Funcionales (Diapositivas 53--54)]
\textbf{A) RNF expresados directamente con plantilla de Historia de Usuario:}
\begin{itemize}[itemsep=2pt]
    \item \textit{``\textbf{Como} comprador, \textbf{quiero} que el sitio web esté disponible el \textbf{98\% de las veces} que intento acceder a él \textbf{para} poder realizar mi compra y no tener que buscar otro lugar para adquirir el producto.''} (Disponibilidad).
    \item \textit{``\textbf{Como} cliente de banca online, \textbf{quiero} que el saldo de mi cuenta corriente se muestre en \textbf{5 segundos} cuando lo solicite \textbf{para} no frustrarme y poder ver si tengo fondos para pagar mis facturas.''} (Rendimiento / Tiempo de respuesta).
\end{itemize}

\textbf{B) RNF expresados como Criterios de Aceptación dentro de una HU funcional:}
\begin{itemize}[itemsep=2pt]
    \item \textbf{HU Funcional:} \textit{``Como analista financiero, quiero ver las transacciones mensuales de mis clientes para poder aconsejarles sobre su salud financiera.''}
    \item \textbf{Criterio RNF 1 (Rendimiento):} \textit{El sistema muestra todas las transacciones que cumplen los parámetros de búsqueda en los \textbf{10 segundos} siguientes a la recepción de la solicitud.}
    \item \textbf{Criterio RNF 2 (Seguridad):} \textit{Las transacciones de los clientes etiquetados como confidenciales sólo se muestran a los usuarios con \textbf{nivel de seguridad 2}.}
\end{itemize}
\end{ejemplo}

\section{¿Cuándo es buena una HU? El Modelo I.N.V.E.S.T. (Diapositivas 55--57)}

Propuesto por \textbf{Bill Wake}, el acrónimo \textbf{I.N.V.E.S.T.} define las seis cualidades que debe reunir toda Historia de Usuario de calidad antes de entrar en una iteración:

\vspace{4pt}
\noindent
\renewcommand{\arraystretch}{1.4}
\begin{tabularx}{0.99\linewidth}{>{\centering\arraybackslash}p{1.2cm} >{\bfseries\sffamily\color{mdaDarkBlue}}p{2.9cm} X}
\toprule
\rowcolor{mdaDarkBlue}
\textbf{\color{white}\normalsize Letra} & \textbf{\color{white}Atributo} & \textbf{\color{white}\textsf{Definición Oficial y Pregunta de Control (Diapositivas 56--57)}} \\
\midrule
\colorbox{mdaDarkBlue}{\makebox[0.65cm]{\color{white}\Large\bfseries I}} & Independiente\newline{\scriptsize(\textit{Independent})} & Una historia debería ser independiente de otras (lo más posible). \textit{¿Puede sostenerse por sí misma sin depender rígidamente de otras historias?} \\
\rowcolor{mdaGray}
\colorbox{mdaDarkBlue}{\makebox[0.65cm]{\color{white}\Large\bfseries N}} & Negociable\newline{\scriptsize(\textit{Negotiable})} & Las HU no son un contrato cerrado ni un requerimiento rígido que el software deba implementar literalmente. \textit{¿Es una propuesta abierta a la discusión colaborativa en lugar de un contrato fijo?} \\
\colorbox{mdaDarkBlue}{\makebox[0.65cm]{\color{white}\Large\bfseries V}} & Valiosa\newline{\scriptsize(\textit{Valuable})} & Tiene que tener valor para el cliente (para el usuario o para el comprador). El cliente las escribe, no son un contrato y son negociables. \textit{¿Aporta un beneficio tangible, demostrable e identificable al usuario o al negocio?} \\
\rowcolor{mdaGray}
\colorbox{mdaDarkBlue}{\makebox[0.65cm]{\color{white}\Large\bfseries E}} & Estimable\newline{\scriptsize(\textit{Estimable})} & Debe poder estimarse para permitir que se pueda priorizar y planificar la historia. \textit{¿Tiene el equipo suficiente información y contexto para calcular su complejidad?} \\
\colorbox{mdaDarkBlue}{\makebox[0.65cm]{\color{white}\Large\bfseries S}} & Suscinta / Pequeña\newline{\scriptsize(\textit{Small})} & Una buena historia debe ser pequeña en esfuerzo (\textbf{2--3 personas/semana de trabajo}, o entre $\mathbf{1/10}$ y $\mathbf{1/6}$ de la velocidad). \textit{¿Es lo suficientemente compacta para completarse con éxito dentro de un solo Sprint?} \\
\rowcolor{mdaGray}
\colorbox{mdaDarkBlue}{\makebox[0.65cm]{\color{white}\Large\bfseries T}} & Testeable\newline{\scriptsize(\textit{Testable})} & Una historia necesita poder probarse objetivamente para que ocurra la etapa de ``Confirmación''. \textit{¿Posee criterios de aceptación claros que permitan validar su correcto funcionamiento?} \\
\bottomrule
\end{tabularx}
\renewcommand{\arraystretch}{1.0}
